% ---------------------------------------------------------------------------
% perturbed_flow.tex  --  fragment for insertion into Appendix A of the manuscript
% (symmetry-4596052, revision 1; answer to Reviewer 1, point (iii)).
% Uses only the MDPI environments Theorem, Lemma, Proposition, Corollary, Remark and
% proof, plus amsmath; no private macros.  The statements of Sections A.1 and the
% spatial step of Section A.2 are machine-checked in PerturbedFlow.lean
% (Lean 4.33.1 + Mathlib); see Table~\ref{tab:lean-perturbed}.
% ---------------------------------------------------------------------------

\section{A-Posteriori Energy Balance for an Arbitrary Trajectory}\label{app:perturbed}

Theorem~\ref{thm:app-gradflow} assumes that the curve $c$ is an \emph{exact} gradient flow, $c'=-\nabla E(c)$.
A trained Spectral-PINN is not: the coefficient trajectory $c(t)$ satisfies the Allen--Cahn equation only up to a \emph{residual}, and the Neumann condition only up to a boundary defect.
This appendix removes the exactness assumption.
It proves, for an \emph{arbitrary} differentiable trajectory, an energy balance whose right-hand side contains only the residual, the boundary defect and the dissipation, and it identifies these quantities for the polynomial states $u(x,t)=\sum_{k=0}^{K}c_k(t)T_k(x)$ produced by the network.
The exact gradient flow is the special case of zero residual and zero defect.
Nothing in this appendix asserts that the trained network has a small residual; the residual and the defect are \emph{a-posteriori} quantities that can be computed for a given network (Remark~\ref{rem:compute}).

\subsection{Abstract Statements}\label{app:perturbed-abstract}

\begin{Lemma}[Algebra of one instant]\label{lem:balance}
Let $H$ be a real inner-product space and let $G,r,w\in H$ satisfy $w=-G+r$. Then
\begin{align}
\langle G,w\rangle &= -\|w\|^2+\langle r,w\rangle, \label{eq:bal1}\\
\langle G,w\rangle &= \tfrac12\|r\|^2-\tfrac12\|w\|^2-\tfrac12\|G\|^2, \label{eq:bal2}\\
\langle G,w\rangle &= \tfrac14\|r\|^2-\bigl\|G-\tfrac12 r\bigr\|^2. \label{eq:bal3}
\end{align}
In particular $\langle G,w\rangle\le\tfrac12\|r\|^2-\tfrac12\|w\|^2$ and $\langle G,w\rangle\le\tfrac14\|r\|^2$.
\end{Lemma}
\begin{proof}
Since $G=r-w$, $\langle G,w\rangle=\langle r,w\rangle-\langle w,w\rangle$, which is \eqref{eq:bal1}.
Since $r=w+G$, $\|r\|^2=\|w\|^2+2\langle w,G\rangle+\|G\|^2$; solving for $\langle G,w\rangle$ gives \eqref{eq:bal2}.
Finally $\bigl\|G-\tfrac12 r\bigr\|^2=\|G\|^2-\langle G,r\rangle+\tfrac14\|r\|^2$, so $\tfrac14\|r\|^2-\|G-\tfrac12 r\|^2=\langle G,r\rangle-\|G\|^2=\langle G,r-G\rangle=\langle G,w\rangle$, which is \eqref{eq:bal3}.
The two inequalities follow by dropping $-\tfrac12\|G\|^2\le0$ in \eqref{eq:bal2} and $-\|G-\tfrac12 r\|^2\le0$ in \eqref{eq:bal3}.
\end{proof}

\begin{Theorem}[A-posteriori energy balance, perturbed gradient flow]\label{thm:perturbed}
Let $F$ be a real Hilbert space and let $E:F\to\mathbb{R}$ be Fr\'echet differentiable at $c(t_0)$ with gradient $\nabla E(c(t_0))$.
Let $c:\mathbb{R}\to F$ be differentiable at $t_0$ and put
\[
r:=c'(t_0)+\nabla E(c(t_0)),\qquad\text{i.e.}\qquad c'(t_0)=-\nabla E(c(t_0))+r .
\]
(Every differentiable curve has such a residual $r$; $r=0$ is the exact gradient flow of Theorem~\ref{thm:app-gradflow}.)
Then $t\mapsto E(c(t))$ is differentiable at $t_0$ and
\begin{align}
\frac{d}{dt}\bigg|_{t_0}E(c(t)) &= -\|c'\|^2+\langle r,c'\rangle \label{eq:pert-id}\\
&=\tfrac12\|r\|^2-\tfrac12\|c'\|^2-\tfrac12\|\nabla E(c)\|^2 \label{eq:pert-exact}\\
&\le \tfrac12\|r\|^2-\tfrac12\|c'\|^2, \label{eq:pert-ineq}
\end{align}
and also $\frac{d}{dt}E(c(t))\le\tfrac14\|r\|^2$ (all quantities at $t_0$). Moreover, if $\|r\|\le\theta\|c'\|$ for some $\theta\ge0$, then $\frac{d}{dt}E(c(t))\le-(1-\theta)\|c'\|^2$; in particular the energy is non-increasing at $t_0$ whenever $\|r\|\le\|c'\|$.
\end{Theorem}
\begin{proof}
By the chain rule, $\frac{d}{dt}E(c(t))\big|_{t_0}=DE(c(t_0))[c'(t_0)]=\langle\nabla E(c(t_0)),c'(t_0)\rangle=\langle G,w\rangle$ with $G:=\nabla E(c(t_0))$ and $w:=c'(t_0)=-G+r$.
Equations \eqref{eq:pert-id}--\eqref{eq:pert-ineq} and the bound $\tfrac14\|r\|^2$ are \eqref{eq:bal1}, \eqref{eq:bal2} and the inequalities of Lemma~\ref{lem:balance}.
If $\|r\|\le\theta\|w\|$, the Cauchy--Schwarz inequality gives $\langle r,w\rangle\le\|r\|\|w\|\le\theta\|w\|^2$, and \eqref{eq:pert-id} yields the last claim.
\end{proof}

\begin{Remark}\label{rem:sharp}
(i) For $r=0$, \eqref{eq:pert-id} is $dE/dt=-\|c'\|^2$, i.e.\ Theorem~\ref{thm:app-gradflow}.
(ii) Inequality \eqref{eq:pert-ineq} keeps the dissipation $-\tfrac12\|c'\|^2$ and costs the constant $\tfrac12$ in front of $\|r\|^2$.
If the dissipation is discarded, the sharp constant is $\tfrac14$: equality in $\frac{d}{dt}E\le\frac14\|r\|^2$ holds iff $\nabla E(c)=\tfrac12 r$, i.e.\ iff $c'=\tfrac12 r=\nabla E(c)$ (the curve climbs the energy at half the residual speed), by \eqref{eq:bal3}.
The constant $\tfrac14$ is optimal for a bound in terms of $\|r\|$ alone: for $F=\mathbb{R}$, $E(x)=x$, $c(t)=t$ one has $\nabla E=1$, $c'=1$, $r=2$ and $dE/dt=1=\tfrac14r^2$.
This is a statement about the best possible bound, not a statement that it is attained by a trained network.
(iii) The Corollary following Theorem~\ref{thm:app-gradflow} instantiates $F=L^2(\Omega)$ for the Ginzburg--Landau functional.
This is formal only: $E[u]=\int(\frac{\varepsilon^2}{2}u_x^2+F(u))\,dx$ is not defined on all of $L^2$ and is not Fr\'echet differentiable there with an $L^2$ gradient unless $u$ has two derivatives and satisfies the Neumann condition.
Theorem~\ref{thm:defect} below, in which $E$ lives on a smaller space $V$ embedded in $H=L^2$, is the rigorous formulation; it also makes the role of the Neumann condition explicit.
\end{Remark}

\begin{Theorem}[Energy balance with a boundary defect]\label{thm:defect}
Let $V$ be a real normed space, $H$ a real inner-product space and $\iota:V\to H$ a continuous linear map (for Allen--Cahn: the coefficient space $V=\mathbb{R}^{K+1}$ or $H^1$, embedded in $H=L^2(-1,1)$).
Let $E:V\to\mathbb{R}$ be Fr\'echet differentiable at $c(t_0)$ with differential $DE(c(t_0))\in V'$, and let $c:\mathbb{R}\to V$ be differentiable at $t_0$.
Assume that there are $G\in H$ and a real number $\beta$ (the value $\beta=\beta(c')$ of a defect functional $\beta\in V'$) such that, at $t_0$,
\begin{equation}\label{eq:defect-split}
DE(c)[c']=\langle G,\iota c'\rangle+\beta
\end{equation}
(for every direction $v$: $DE(c)[v]=\langle G,\iota v\rangle+\beta(v)$), and put $r:=\iota c'+G\in H$, i.e.\ $\iota c'=-G+r$.
Then $t\mapsto E(c(t))$ is differentiable at $t_0$ and
\begin{equation}\label{eq:defect-id}
\begin{aligned}
\frac{d}{dt}E(c(t)) &= -\|\iota c'\|^2+\langle r,\iota c'\rangle+\beta\\
&=\tfrac12\|r\|^2-\tfrac12\|\iota c'\|^2-\tfrac12\|G\|^2+\beta\\
&\le\tfrac12\|r\|^2-\tfrac12\|\iota c'\|^2+\beta,
\end{aligned}
\end{equation}
and also $\frac{d}{dt}E(c(t))\le\tfrac14\|r\|^2+\beta$.
\end{Theorem}
\begin{proof}
The chain rule gives $\frac{d}{dt}E(c(t))=DE(c)[c']$, which equals $\langle G,\iota c'\rangle+\beta$ by \eqref{eq:defect-split}.
Apply Lemma~\ref{lem:balance} in $H$ with $w=\iota c'$.
\end{proof}

\begin{Corollary}[Integrated bounds]\label{cor:integrated}
Let $c$ and $E$ be as in Theorem~\ref{thm:defect}, but now on a time interval: assume that $c$ is defined and differentiable on an open interval containing $[0,T]$, that $E$ is differentiable at every $c(t)$, and that \eqref{eq:defect-split} holds at every $t\in[0,T]$ with $G(t)$, $\beta(t)$ and residual $r(t)=\iota c'(t)+G(t)$.
Suppose $\|r(t)\|\le\delta$ and $|\beta(t)|\le b$ for all $t\in[0,T]$. Then for $0\le s\le t\le T$,
\begin{equation}\label{eq:integrated}
E(c(t))\le E(c(s))+\Bigl(\frac{\delta^2}{2}+b\Bigr)(t-s),
\end{equation}
and the constant $\delta^2/2$ may be replaced by the sharp constant $\delta^2/4$.
In particular (no defect, $V=H=F$, $\iota=\mathrm{id}$, $G=\nabla E(c)$): if $c'=-\nabla E(c)+r$ with $\|r(t)\|\le\delta$, then $E(c(t))\le E(c(0))+\frac{\delta^2}{2}\,t$; for $\delta=0$ the energy is non-increasing on $[0,T]$.

If moreover $t\mapsto DE(c(t))[c'(t)]$, $\|r(t)\|^2$, $\|\iota c'(t)\|^2$ and $\beta(t)$ are continuous on $[0,T]$, then
\begin{equation}\label{eq:integrated-L2}
E(c(t))+\frac12\int_0^t\|\iota c'(s)\|^2\,ds\;\le\;E(c(0))+\frac12\int_0^t\|r(s)\|^2\,ds+\int_0^t\beta(s)\,ds,\qquad t\in[0,T].
\end{equation}
\end{Corollary}
\begin{proof}
By the chain rule $\psi(t):=E(c(t))$ is differentiable on $[0,T]$, hence continuous there, and by Theorem~\ref{thm:defect} $\psi'(\xi)\le\tfrac12\|r(\xi)\|^2+\beta(\xi)\le\tfrac12\delta^2+b$ for every $\xi\in(0,T)$ (the term $-\tfrac12\|\iota c'\|^2\le0$ is dropped).
The mean-value theorem on $[s,t]$ gives $\psi(t)-\psi(s)=\psi'(\xi)(t-s)\le(\tfrac12\delta^2+b)(t-s)$.
The same argument with $\psi'\le\tfrac14\|r\|^2+\beta$ gives the constant $\delta^2/4$.
For \eqref{eq:integrated-L2}, the continuity assumptions make $\psi'(s)=DE(c(s))[c'(s)]$ continuous, so by the fundamental theorem of calculus $\psi(t)-\psi(0)=\int_0^t\psi'(s)\,ds$, and integrating $\psi'\le\frac12\|r\|^2-\frac12\|\iota c'\|^2+\beta$ over $[0,t]$ gives the claim.
\end{proof}

\begin{Corollary}[Discrete-time form]\label{cor:discrete}
Let $E_0,E_1,\dots,E_N$ be real numbers and let $r_n,d_n$ ($n<N$) be vectors of a normed space such that $E_{n+1}-E_n\le\tfrac12\|r_n\|^2-\tfrac12\|d_n\|^2$ for $n=0,\dots,N-1$. Then
\[
E_N+\tfrac12\sum_{n=0}^{N-1}\|d_n\|^2\le E_0+\tfrac12\sum_{n=0}^{N-1}\|r_n\|^2,
\qquad\text{hence}\qquad E_N\le E_0+\tfrac12\sum_{n=0}^{N-1}\|r_n\|^2 ,
\]
and $E_N\le E_0+\frac N2\delta^2$ if $\|r_n\|\le\delta$ for all $n$.
With step weights $\tau_n\ge0$ and $E_{n+1}-E_n\le\tau_n\bigl(\tfrac12\|r_n\|^2-\tfrac12\|d_n\|^2\bigr)$ one obtains $E_N+\tfrac12\sum\tau_n\|d_n\|^2\le E_0+\tfrac12\sum\tau_n\|r_n\|^2$ and $E_N\le E_0+\frac{\delta^2}{2}\sum_{n<N}\tau_n$.
\end{Corollary}
\begin{proof}
Sum the hypothesis over $n=0,\dots,N-1$; the left-hand sides telescope to $E_N-E_0$.
The second and third claims drop the non-negative term $\frac12\sum\|d_n\|^2$ and bound $\|r_n\|^2\le\delta^2$.
The weighted statement is the same computation with the weights $\tau_n$.
\end{proof}

\begin{Remark}[Where the discrete hypothesis comes from]
For a time-stepping scheme the hypothesis of Corollary~\ref{cor:discrete} is a one-step energy inequality.
For the checkpoint energies $E_n=E(c(t_n))$ of a trained network in the defect-free case it holds with the non-negative numbers $R_n,D_n$ defined by $R_n^2:=\int_{t_n}^{t_{n+1}}\|r\|^2\,dt$ and $D_n^2:=\int_{t_n}^{t_{n+1}}\|\iota c'\|^2\,dt$ in place of $\|r_n\|,\|d_n\|$, by integrating Theorem~\ref{thm:perturbed} over $[t_n,t_{n+1}]$ (the energy of a polynomial state is a polynomial in the coefficients, hence $t\mapsto E(c(t))$ is $C^1$ and the fundamental theorem of calculus applies).
With a boundary defect, the term $\int_{t_n}^{t_{n+1}}\beta\,dt$ has to be added on the right.
That integration step is not part of the machine-checked file; Corollary~\ref{cor:discrete} is checked as a statement about arbitrary real sequences.
\end{Remark}

\subsection{Link to the Allen--Cahn Equation}\label{app:perturbed-pde}

Let $\varepsilon>0$, $f(u)=u^3-u=F'(u)$, $F(u)=\tfrac14(u^2-1)^2$, $L^2=L^2(-1,1)$ with $\langle v,w\rangle=\int_{-1}^{1}vw\,dx$ and $\|\cdot\|=\|\cdot\|_{L^2}$.
Let $I\supset[0,T]$ be an open interval, $c=(c_0,\dots,c_K)\in C^1(I;\mathbb{R}^{K+1})$, and
\[
u(x,t)=\sum_{k=0}^{K}c_k(t)\,T_k(x),
\]
where $T_k$ are the Chebyshev polynomials (the argument uses only that the $T_k$ are fixed real polynomials; no property of the network is used beyond $c\in C^1$).
Define the PDE residual and the energy of the learned trajectory,
\[
r(x,t):=u_t-\varepsilon^2u_{xx}+f(u),\qquad E(t):=\int_{-1}^{1}\Bigl(\frac{\varepsilon^2}{2}u_x^2+F(u)\Bigr)dx .
\]
Note that $r(\cdot,t)$ is a polynomial in $x$ of degree at most $3K$, so $r(\cdot,t)\in L^2$ and all norms below are finite.

\begin{Proposition}[Energy identity for polynomial-in-$x$ states]\label{prop:pde-identity}
$E\in C^1(I)$ and, for every $t\in I$,
\begin{equation}\label{eq:pde-identity}
\begin{aligned}
\frac{dE}{dt}&=-\|u_t\|^2+\langle r,u_t\rangle+\varepsilon^2\bigl[u_xu_t\bigr]_{x=-1}^{x=1}\\
&=\tfrac12\|r\|^2-\tfrac12\|u_t\|^2-\tfrac12\bigl\|{-\varepsilon^2u_{xx}}+f(u)\bigr\|^2+\varepsilon^2\bigl[u_xu_t\bigr]_{x=-1}^{x=1},
\end{aligned}
\end{equation}
where $[g]_{x=-1}^{x=1}=g(1,t)-g(-1,t)$.
\end{Proposition}
\begin{proof}
\emph{Step 1: $E$ is a polynomial in the coefficients; this is why differentiation under the integral sign is legitimate.}
For $a=(a_0,\dots,a_K)\in\mathbb{R}^{K+1}$ put $u_a(x)=\sum_ka_kT_k(x)$ and
\[
\mathcal{E}(a):=\int_{-1}^{1}\Bigl(\frac{\varepsilon^2}{2}(u_a')^2+\tfrac14u_a^4-\tfrac12u_a^2+\tfrac14\Bigr)dx .
\]
Expanding the powers, $(u_a')^2=\sum_{j,k}a_ja_kT_j'T_k'$ and $u_a^m=\sum_{k_1,\dots,k_m}a_{k_1}\cdots a_{k_m}T_{k_1}\cdots T_{k_m}$; hence the integrand is a finite sum $\sum_\alpha a^\alpha q_\alpha(x)$ with $|\alpha|\le4$ and real polynomials $q_\alpha$, and
$\mathcal{E}(a)=\sum_\alpha a^\alpha\int_{-1}^{1}q_\alpha\,dx$ is a real polynomial in $a$.
Thus $\mathcal{E}\in C^\infty(\mathbb{R}^{K+1})$, and since $c\in C^1$ and $E(t)=\mathcal{E}(c(t))$, the chain rule gives $E\in C^1(I)$ and $E'(t)=D\mathcal{E}(c(t))[c'(t)]$.

\emph{Step 2: the differential of $\mathcal{E}$.}
Fix $a,v\in\mathbb{R}^{K+1}$, write $w:=u_v$, so that $u_{a+sv}=u_a+sw$.
The integrand of $\mathcal{E}(a+sv)$ is $\frac{\varepsilon^2}{2}(u_a'+sw')^2+F(u_a+sw)$, a polynomial of degree at most $4$ in $s$ whose coefficients are polynomials in $x$.
Integration over $[-1,1]$ commutes with finite sums, so $s\mapsto\mathcal{E}(a+sv)$ is a polynomial in $s$ whose coefficient of $s$ is the integral of the coefficient of $s$ in the integrand, namely $\varepsilon^2u_a'w'+F'(u_a)w$. Hence
\begin{equation}\label{eq:DE}
D\mathcal{E}(a)[v]=\int_{-1}^{1}\bigl(\varepsilon^2u_a'w'+f(u_a)\,w\bigr)dx .
\end{equation}
No limit of integrals is exchanged with a derivative except for a finite sum, so no dominated-convergence argument is needed.

\emph{Step 3: integration by parts in $x$.}
Let $p,q$ be real polynomials. Then $(pq)'=p'q+pq'$, and by the fundamental theorem of calculus $\int_{-1}^{1}(pq)'dx=p(1)q(1)-p(-1)q(-1)$; hence $\int_{-1}^{1}p\,q'\,dx=[pq]_{-1}^{1}-\int_{-1}^{1}p'q\,dx$.
Apply this with $p=u_a'$ and $q=w$ in \eqref{eq:DE}:
\begin{equation}\label{eq:DE-split}
D\mathcal{E}(a)[v]=\underbrace{\int_{-1}^{1}\bigl(-\varepsilon^2u_a''+f(u_a)\bigr)\,w\,dx}_{=\ \langle G_a,\,\iota v\rangle}\;+\;\underbrace{\varepsilon^2\bigl[u_a'\,w\bigr]_{x=-1}^{x=1}}_{=\ \beta_a(v)},
\end{equation}
with $G_a:=-\varepsilon^2u_a''+f(u_a)\in L^2$ and $\iota v:=u_v=w$.
This is exactly the splitting \eqref{eq:defect-split}, with the defect functional $\beta_a(v)=\varepsilon^2[u_a'u_v]_{-1}^{1}$; it vanishes for all $v$ iff $u_a'(\pm1)=0$, i.e.\ iff the Neumann condition holds.

\emph{Step 4: insert the trajectory.}
Take $a=c(t)$ and $v=c'(t)$. Since the $T_k$ do not depend on $t$ and $c_k\in C^1$, $u_t=\sum_kc_k'T_k=u_{c'}$ and $u_{xt}=\partial_x(u_t)$.
Thus $E'(t)=\langle G,u_t\rangle+\varepsilon^2[u_xu_t]_{-1}^{1}$ with $G=-\varepsilon^2u_{xx}+f(u)$.
By the definition of $r$, $u_t=-G+r$ in $L^2$, so Lemma~\ref{lem:balance} with $w=u_t$ gives $\langle G,u_t\rangle=-\|u_t\|^2+\langle r,u_t\rangle=\frac12\|r\|^2-\frac12\|u_t\|^2-\frac12\|G\|^2$, which is \eqref{eq:pde-identity}.
\end{proof}

% Numerical sanity check of the signs/forms in (eq:pde-identity): see check_identity.py (same folder).
(As a sanity check of signs and constants, \eqref{eq:pde-identity} was compared with a central finite difference of $E(t)$, evaluated with exact Gauss--Legendre quadrature, for a random smooth coefficient trajectory with $K=7$ and $\varepsilon=0.3$ at three times: the relative difference was below $10^{-9}$, which is the accuracy of the finite difference. This is a plausibility check, not part of the proof.)

\begin{Corollary}[The trained trajectory has controlled energy growth]\label{cor:spinn-bound}
Let $u$ be as above and $0\le s\le t\le T$. Put
\[
\delta:=\max_{\tau\in[0,T]}\|r(\cdot,\tau)\|_{L^2},\qquad
b:=\varepsilon^2\max_{\tau\in[0,T]}\Bigl(|u_x(1,\tau)\,u_t(1,\tau)|+|u_x(-1,\tau)\,u_t(-1,\tau)|\Bigr).
\]
Then
\[
E(t)\le E(s)+\Bigl(\frac{\delta^2}{2}+b\Bigr)(t-s),
\]
and, with $\beta(\tau):=\varepsilon^2[u_xu_t]_{x=-1}^{x=1}(\tau)$,
\[
E(t)+\frac12\int_0^t\|u_t\|^2\,d\tau\le E(0)+\frac12\int_0^t\|r\|^2\,d\tau+\int_0^t\beta\,d\tau .
\]
If the Neumann condition holds exactly, $u_x(\pm1,\tau)=0$, then $b=0$ and $\beta\equiv0$, and the bound reduces to $E(t)\le E(s)+\frac{\delta^2}{2}(t-s)$; if moreover the PDE holds exactly ($\delta=0$) one recovers $dE/dt=-\|u_t\|^2\le0$ (Theorem~\ref{thm:energy}).
\end{Corollary}
\begin{proof}
Apply Corollary~\ref{cor:integrated} with $V=\mathbb{R}^{K+1}$, $H=L^2$, $\iota c=u_c$, $E=\mathcal{E}$, $G(t)=-\varepsilon^2u_{xx}+f(u)$ and $\beta(t)=\varepsilon^2[u_xu_t]_{-1}^{1}$; the splitting \eqref{eq:defect-split} holds for every $t$ by \eqref{eq:DE-split}.
The maxima exist and the continuity hypotheses of \eqref{eq:integrated-L2} hold because the coefficients (in $x$) of $r$, $u_t$, $u_x$ and of the integrand of $D\mathcal{E}(c)[c']$ are polynomials in the continuous functions $c_k,c_k'$ ($r$ contains $u_t=\sum_kc_k'T_k$ but no second time derivative), so $\tau\mapsto\|r(\cdot,\tau)\|^2$, $\|u_t(\cdot,\tau)\|^2$, $\beta(\tau)$ and $D\mathcal{E}(c(\tau))[c'(\tau)]$ are continuous on the compact interval $[0,T]$ when $c\in C^1$.
\end{proof}

\begin{Remark}[What this does and does not say about the trained network]\label{rem:compute}
(i) The theorem is conditional: it converts a bound on the PDE residual and on the Neumann defect into a bound on energy growth. It does not prove that training makes $\delta$ or $b$ small, and the resulting bound $\frac{\delta^2}{2}T$ can be far larger than the actual energy variation.
The values of $\delta$ and $b$ must be reported for each trained network before any quantitative conclusion is drawn; they are computable to rounding error, because $r(\cdot,t)$ is a polynomial of degree $\le3K$ in $x$, so $\|r(\cdot,t)\|^2$ is integrated exactly by Gauss--Legendre quadrature with $N\ge3K+1$ nodes, and $u_x(\pm1,t)$, $u_t(\pm1,t)$ are evaluated directly from the coefficients.
(ii) The training objective contains the mean-square residual at collocation points and a boundary penalty on $u_x(\pm1,t)$.
In the code accompanying the submission (\texttt{spectral\_pinn.py}, default settings) the residual is evaluated on $100$ uniformly spaced times and $300$ uniformly spaced points of $[-1,1]$, so the PDE loss is a mean over this grid, which approximates $\frac{1}{2T}\int_0^T\|r\|^2\,dt$ (the factor $2=|\Omega|$) but is not an exact quadrature of the degree-$6K$ polynomial $r^2$; and the boundary term penalizes $u_x(\pm1,t)/K^2$ (the boundary rows are rescaled by $K^2$), i.e.\ $u_x(\pm1,t)^2$ with weight $K^{-4}\lambda_{\mathrm{bc}}$, which is a very weak control of $b$ for large $K$.
Hence the loss values reported during training are proxies for $\delta$ and $b$ and must not be substituted for them; both should be evaluated on the trained network as described in (i).
With the integral form \eqref{eq:integrated-L2}, the energy increase is controlled by the time-integrated squared residual and the boundary term, which are the two quantities that the PDE and boundary losses are designed to make small.
(iii) The bound does not involve the energy-penalty term of the loss.
This is consistent with the observation that the energy curves with and without the penalty are nearly indistinguishable, but it does not prove why: it only shows that a small residual and a small Neumann defect already constrain the energy growth.
(iv) The initial condition $u_0(x)=x^2\cos(\pi x)$ of the submitted version has $u_{0,x}(-1)=2\neq0$, so for a network that fits $u_0$ the factor $u_x(-1,0)$ in the defect term is not small at $t=0$; whether $b$ is small then depends on $u_t(\pm1,t)$ and has to be measured, not assumed. With a Neumann-compatible initial condition this particular source of a non-zero defect at $t=0$ disappears.
\end{Remark}

\subsection{Machine Verification}\label{app:perturbed-lean}

All statements of Section~\ref{app:perturbed-abstract} and the spatial integration by parts of Step~3 above were machine-checked in Lean~4 with Mathlib (file \texttt{PerturbedFlow.lean}); the file compiles against Lean~4.33.1 and the corresponding Mathlib release with no errors, no warnings and no unproven placeholders, and each theorem depends only on the axioms \texttt{propext}, \texttt{Classical.choice} and \texttt{Quot.sound} (checked with \texttt{\#print axioms}).
The abstract statements are proved for general real Hilbert (or normed) spaces by Mathlib's chain rule (\texttt{HasFDerivAt.comp\_hasDerivAt\_of\_eq}) and mean-value theorem (\texttt{Convex.image\_sub\_le\_mul\_sub\_of\_deriv\_le}), the integral form by the fundamental theorem of calculus, and the integration by parts for polynomials by \texttt{intervalIntegral.}\allowbreak\texttt{integral\_mul\_deriv\_eq\_deriv\_mul}.

\begin{table}[H]
\caption{Correspondence between the statements of this appendix and the Lean theorems of \texttt{PerturbedFlow.lean}.\label{tab:lean-perturbed}}
\small
\begin{tabularx}{\textwidth}{l>{\raggedright\arraybackslash}X}
\toprule
\textbf{Statement} & \textbf{Lean theorem(s)}\\
\midrule
Lemma~\ref{lem:balance} & \texttt{inner\_balance}, \texttt{inner\_balance\_exact}, \texttt{inner\_balance\_quarter}, \texttt{inner\_balance\_le}, \texttt{inner\_balance\_le\_quarter}\\
Theorem~\ref{thm:perturbed}, \eqref{eq:pert-id} & \texttt{perturbedFlow\_hasDerivAt\_energy}, \texttt{perturbedFlow\_hasDerivAt\_energy\_gradient}\\
Theorem~\ref{thm:perturbed}, \eqref{eq:pert-exact} & \texttt{perturbedFlow\_hasDerivAt\_energy\_exact}\\
Theorem~\ref{thm:perturbed}, \eqref{eq:pert-ineq} & \texttt{perturbedFlow\_deriv\_energy\_le}, \texttt{perturbedFlow\_deriv\_energy\_le\_deriv}\\
Theorem~\ref{thm:perturbed}, $\|r\|\le\theta\|c'\|$ & \texttt{perturbedFlow\_deriv\_energy\_le\_of\_relative}, \texttt{perturbedFlow\_deriv\_energy\_nonpos\_of\_residual\_le}\\
Remark~\ref{rem:sharp}(i) ($r=0$) & \texttt{perturbedFlow\_exact\_case}\\
Theorem~\ref{thm:defect}, identity & \texttt{defectFlow\_hasDerivAt\_energy}, \texttt{defectFlow\_hasDerivAt\_energy\_functional}, \texttt{defectFlow\_hasDerivAt\_energy\_exact}\\
Theorem~\ref{thm:defect}, inequalities & \texttt{defectFlow\_deriv\_energy\_le}, \texttt{defectFlow\_deriv\_energy\_le\_quarter}\\
Corollary~\ref{cor:integrated}, \eqref{eq:integrated} & \texttt{energy\_le\_of\_deriv\_le}, \texttt{defectFlow\_energy\_growth\_bound}, \texttt{defectFlow\_energy\_bound}, \texttt{defectFlow\_energy\_bound\_quarter}, \texttt{perturbedFlow\_energy\_growth\_bound}, \texttt{perturbedFlow\_energy\_bound}, \texttt{perturbedFlow\_energy\_bound\_quarter}, \texttt{gradientFlow\_energy\_le\_initial}\\
Corollary~\ref{cor:integrated}, \eqref{eq:integrated-L2} & \texttt{energy\_integral\_le}, \texttt{defectFlow\_energy\_integral\_bound}, \texttt{defectFree\_energy\_integral\_bound}\\
Corollary~\ref{cor:discrete} & \texttt{energy\_le\_of\_increment\_le}, \texttt{discrete\_energy\_balance}, \texttt{discrete\_energy\_bound}, \texttt{discrete\_energy\_bound\_uniform}, \texttt{discrete\_energy\_balance\_weighted}, \texttt{discrete\_energy\_bound\_weighted}\\
Proposition~\ref{prop:pde-identity}, Step~3 and \eqref{eq:DE-split} & \texttt{ibp\_polynomial}, \texttt{first\_variation\_split}\\
\bottomrule
\end{tabularx}
\end{table}

Not machine-checked are: Step~1 of Proposition~\ref{prop:pde-identity} (that $\mathcal{E}$ is a polynomial in the coefficients), the identification of the abstract objects $E,G,\beta,r$ with the Allen--Cahn quantities in Corollary~\ref{cor:spinn-bound} (this identification is the content of Steps~2--4 and uses only the verified integration by parts), the fundamental-theorem-of-calculus step in the remark following Corollary~\ref{cor:discrete}, and, of course, the empirical values of $\delta$ and $b$ for the trained networks.
